% Activity diagram sistem Rupiah Digital.
% Ketujuh gambar di-render dari blok bernama pada berkas PlantUML:
%   bi-coin-fabric/docs/diagrams/plantuml/activity-diagrams.puml
% Untuk merender:
%   cd /path/to/bi-coin && rtk python render-activity-diagrams.py
% Di-input dari BAB III §3.2.4.

\newcommand{\actfig}[3]{%
  \begin{figure}[H]\centering
  \IfFileExists{contents/figures/activity/#1.png}{%
    \includegraphics[width=\textwidth,height=0.78\textheight,keepaspectratio]%
      {contents/figures/activity/#1.png}%
  }{%
    \fbox{\parbox{0.9\textwidth}{\centering\small
      \textit{[Gambar belum di-render. Jalankan:
      \texttt{rtk python render-activity-diagrams.py}]}\\
      \texttt{#1.png}}}%
  }%
  \caption{#2}
  \label{#3}
  \end{figure}
}

\newcommand{\actfiglandscape}[3]{%
  \begin{figure}[H]\centering
  \IfFileExists{contents/figures/activity/#1.png}{%
    \rotatebox{90}{\includegraphics[width=0.72\textheight]%
      {contents/figures/activity/#1.png}}%
  }{%
    \fbox{\parbox{0.9\textwidth}{\centering\small
      \textit{[Gambar belum di-render. Jalankan:
      \texttt{rtk python render-activity-diagrams.py}]}\\
      \texttt{#1.png}}}%
  }%
  \caption{#2}
  \label{#3}
  \end{figure}
}

\actfiglandscape{AD01-Pendaftaran-Peserta}%
  {Activity diagram pendaftaran peserta}{fig:ad01}

\actfig{AD02-Penerbitan-Digital-Rupiah}%
  {Activity diagram penerbitan Digital Rupiah}{fig:ad02}

\actfig{AD03-Distribusi-Likuiditas-PJP}%
  {Activity diagram distribusi likuiditas ke PJP}{fig:ad03}

\actfig{AD05-Transfer-Saldo-Ritel}%
  {Activity diagram transfer saldo ritel}{fig:ad05}


\actfig{AD08-Pembekuan-Peserta}%
  {Activity diagram pembekuan dan pencairan peserta oleh Bank Indonesia}{fig:ad08}

\actfiglandscape{AD10-Laporan-Supervisi}%
  {Activity diagram laporan dan pengawasan supervisi}{fig:ad10}
